\documentclass[./examen.tex]{subfiles}

\begin{document}
\begin{center}
\makebox[\textwidth]{Nombre y Apellido:\enspace\hrulefill}
\end{center}
\begin{questions}
 \question La máquina térmica con eficiencia de $\eta_{\%}=87\%$ produce $W=4186J$ de trabajo. Encontrar: \begin{parts}
        \part[2] La cantidad de calor absorbida del depósito térmico.
        \part[2] La cantidad de calor desechado al depósito térmico.
    \end{parts} 
    
\question Para el funcionamiento de una máquina, se pierde el 30\% del calor entregado, si la fuente caliente está a $T_1=\SI{400}{\celsius}$ y se realiza un trabajo de $W200J$ Calcular:
      \begin{parts}
          \part[1] El calor perdido $Q_2$. 
          \part[1] El rendimiento $\eta$.
          \part[1] El calor entregad $Q_1$.
          \part[2] calcular $T_1$ y $T_2$
          \part[1] Si se quiere duplicar el trabajo realizado ¿Cuál debería ser la temperatura de la fuente caliente $T_1$? 
      \end{parts}
\end{questions}





 \end{document}